% !TEX options=--shell-escape

\documentclass{article}

\usepackage{multicol} % Enables multicol env.
\usepackage{graphicx} % Enables including graphics.
\usepackage[export]{adjustbox} % Allows graphics to have a `center` property.
\usepackage[pdf]{graphviz}
\usepackage{mdframed}
\usepackage{parskip}
\usepackage{amssymb}
\usepackage{titlesec} % Enables custom title sections.
\usepackage{float}


\newmdtheoremenv{theo}{Theorem}
\newmdtheoremenv{definition}{Definition}
\newmdtheoremenv{lemma}{Lemma}

\renewcommand\thesubsection{(\alph{subsection})}
\titleformat{\subsection}
{\normalfont\normalsize}{\thesubsection}{1em}{}

\renewcommand\thesubsubsection{\hspace{10px}\roman{subsubsection}.}
\titleformat{\subsubsection}
{\normalfont\normalsize}{\thesubsubsection}{1em}{}

\newcommand{\YearPath}{../../../LatexConfig} % Get year level path (i.e. Senior)
\newcommand{\SemesterPath}{../../LatexConfig} % Get semester level path (i.e. Semester 1).
\newcommand{\ClassPath}{../LatexConfig} % Get class level path (i.e. CS101).
\newcommand{\AssignmentTitle}{Homework 3}
\newcommand{\AssignmentSub}{~ ~}

\begin{document}

\newcommand{\Professor}{Dr. Thorne Wolf}
\newcommand{\Course}{CS 101.1}
\newcommand{\Professor}{Dr. Thorne Wolf}
\newcommand{\Course}{CS 101.1}
\begin{titlepage}
  \begin{center}
      \vspace*{1cm}

      \Large
      \textbf{\AssignmentTitle}

      \vspace{0.5cm}
      \large
      \AssignmentSub \\
      ~\\
      \normalsize \today

      \vspace{1.5cm}

      \large
      By: \\
      ~\\
      \normalsize
      \vspace{1ex}
      \ifdefined\StudentNames \StudentNames \else Thorne Wolfenbarger\fi
      \vfill

      \vspace{0.8cm}

      \large
      Submitted to: \\
      \Professor \\
      In Partial Fulfillment of the Requirements of \\
      \Course~-~\Semester\\

      \vspace{0.8cm}

      \includegraphics[width=0.2\textwidth,center]{\YearPath/univ}
      ~\\
      College of Engineering\\
      Embry-Riddle Aeronautical University\\
      Prescott, AZ\\

  \end{center}
\end{titlepage}


\section{Question 1}

\subsection{Provide the full statement of at least one theorem, as it’s stated on the presenter’s slides.}
Every $C^*$-algebra looks like a closed $*$-subalgebra of $B(H)$ for some Hilbert space $H$.

\subsection{In the theorem statement above, what type of a proposition is it? Is it conditional? Does it have suppressed quantifiers? If so, rewrite the theorem statement showing these.}
It is a universal statement because the theorem states that ``Every $C^*$-algebra...'' an example of this is:

\begin{center}
  $c(x) = \{f : x \rightarrow \mathbb{C} : $f is continuous$\}$ where x is compact, Hausdorff...
\end{center}

I suppose that the theorem could be conditional because it has to be a $C^*$-algebra.

We can write this as
\begin{center}
  $\forall C^* algebra~\exists H ~s.t. ~C^* algebra~ = ~closed~*$-subalgebra$~ \subseteq B(H)$.
\end{center}

\subsection{Write down a math term/word/notation that seems familiar to you (look out for words like injective, surjective, or notation like function notation).}
Let $(A(\mathbb{D}), \alpha)$ be a dynamical system. Then there exists a Mobius transformation of the form

\begin{center}
  $\tau(w) = \lambda (\frac{w-\mu}{1-\mu\bar w})$
\end{center}
such that $\alpha_1(f)(w) = f(\tau(w))$.

Within this I understand what it means for a system to be dynamical and for a transformation between spaces to exist.

\subsection{What do you understand this word/notation to mean?}
I understand this exerpt to mean there is some transform called a Mobius transform that exists as a function that accepts an argument $w$ and uses it to convert into something that is possibly called Mobius space.

\subsection{Using a resource (wikipedia is fine), provide a formal definition of the word or symbol as it was intended in the presentation.}
Here I look into the definition of a Mobius transformation.

In geometry and complex analysis, a Möbius transformation of the complex plane is a rational function of the form

\begin{center}
  $f(z) = \frac{az+b}{cz+d}$
\end{center}
of one complex variable $z$; here the coefficients $a,b,c,d$ are complex numbers satisfying $ab-bc \neq 0$.

"A Mobius transformation is always a bijective holomorphic function from the Riemann sphere to the Riemann sphere."

"These transformations preserve angles, map every straight line to a line or circle, and map every circle to a line or circle."

All information is probided from (Wikipedia, 2020). These exerpts provide the definition of a Mobius transformation as well as an invarient within the transformation.

\subsection{Write down a math term/word/notation that was totally unfamiliar to you.}
Irrational Rotation.

\subsection{Using a resource (wikipedia is fine), provide a formal definition of the word or symbol as it was intended in the presentation.}
In the mathematical theory of dynamicval systems, an irrational rotation is a map where $\theta$ is an irrational number. Under the identification of a circle with $R/z$ or with the interval $[0,1]$ with the boundary points glued together, this map becomes a rotation of a circle by a proportion $\theta$ of a full revolution.

\subsection{Ask a question that the presentation raised for you that you would like to know more about.}
Why isn't any orbit periodic? Within the same slide, Dr. Hamidi makes a statement "No orbit is periodic." What is a non-periodic orbit?

\newpage

\section{Question 2}
Let $A = {1,3}$ and let $B = {1,2,3,4}$.
\subsection{Describe the set $A \times B$ using the roster method.}
$A \times B = \{(1,1), (1, 2), (1, 3), (1, 4), (3,1), (3, 2), (3, 3), (3, 4)\}$

\subsection{Describe the set $A \times B$ using the roster method.}
$B \times A = \{(1, 1), (1, 3), (2, 1), (2, 3), (3, 1), (3, 3), (4, 1), (4, 3), \}$

\subsection{Are the sets $A \times B$ and $B \times A$ equal? Why or why not?}
They are not equal sets. The cartesian product operator is not a communative operator. It can as the elements within the cartesian product of $A, B$ are ordered pairs, the order of $A, B$ within the operation matters. 

\subsection{Find $A \cup B$}
$A \cup B = \{1, 2, 3, 4\}$

\subsection{Find $A \cap B$}
$A \cap B = \{1, 3\}$

\subsection{Prove the following: Let $A$ and $B$ be two sets. Then $A \subseteq B$ if and only if $A \cup B = B$.}
To say that $A \subseteq B$ we are saying $\forall x \in A (x \in B)$. In set builder notation we can represent the set $A$ as $\{x : x \in A\}$. Since we know $\forall x \in A (x \in B)$ we can additionally constrain the previous set producing $\{x : x \in A\} = \{x : x \in A \land x \in B\}$. These statements are all equivalent, so we can represent this as $\{x : x \in A\} = \{x : x \in A \land x \in B\} \iff \forall x \in A (x \in B)$

Representing $A \cup B$ in set notation provides us with $\{x : x \in A \lor x \in B\}$. If $A \cup B = B$, that means $\{x : x \in A \lor x \in B\} = \{x : x \in B\}$ The only way this could be possible was if $\forall x \in A (x \in B)$ was true. That is, $\forall x \in A (x \in B) \iff \{x : x \in A \lor x \in B\} = \{x : x \in B\}$

We can see that both statements are equivalent to $\forall x \in A (x \in B)$.

\begin{center}
\begin{tabular}{c}
$\{x : x \in A\} = \{x : x \in A \land x \in B\} \iff \forall x \in A (x \in B)$ \\
$\forall x \in A (x \in B) \iff \{x : x \in A \lor x \in B\} = \{x : x \in B\}$\\
$\{x : x \in A\} = \{x : x \in A \land x \in B\} \iff \{x : x \in A \lor x \in B\} = \{x : x \in B\}$ \\
${A \subseteq B \iff A \cup B = B}$
\end{tabular}
\end{center}
Q.E.D
\newpage


\section{Question 3}

\subsection{Provide the full statement of at least one theorem, as it’s stated on the presenter’s slides.}
Every $C^*$-algebra looks like a closed $*$-subalgebra of $B(H)$ for some Hilbert space $H$.

\subsection{In the theorem statement above, what type of a proposition is it? Is it conditional? Does it have suppressed quantifiers? If so, rewrite the theorem statement showing these.}
It is a universal statement because the theorem states that ``Every $C^*$-algebra...'' an example of this is:

\begin{center}
  $c(x) = \{f : x \rightarrow \mathbb{C} : $f is continuous$\}$ where x is compact, Hausdorff...
\end{center}

I suppose that the theorem could be conditional because it has to be a $C^*$-algebra.

We can write this as
\begin{center}
  $\forall C^* algebra~\exists H ~s.t. ~C^* algebra~ = ~closed~*$-subalgebra$~ \subseteq B(H)$.
\end{center}

\subsection{Write down a math term/word/notation that seems familiar to you (look out for words like injective, surjective, or notation like function notation).}
Let $(A(\mathbb{D}), \alpha)$ be a dynamical system. Then there exists a Mobius transformation of the form

\begin{center}
  $\tau(w) = \lambda (\frac{w-\mu}{1-\mu\bar w})$
\end{center}
such that $\alpha_1(f)(w) = f(\tau(w))$.

Within this I understand what it means for a system to be dynamical and for a transformation between spaces to exist.

\subsection{What do you understand this word/notation to mean?}
I understand this exerpt to mean there is some transform called a Mobius transform that exists as a function that accepts an argument $w$ and uses it to convert into something that is possibly called Mobius space.

\subsection{Using a resource (wikipedia is fine), provide a formal definition of the word or symbol as it was intended in the presentation.}
Here I look into the definition of a Mobius transformation.

In geometry and complex analysis, a Möbius transformation of the complex plane is a rational function of the form

\begin{center}
  $f(z) = \frac{az+b}{cz+d}$
\end{center}
of one complex variable $z$; here the coefficients $a,b,c,d$ are complex numbers satisfying $ab-bc \neq 0$.

"A Mobius transformation is always a bijective holomorphic function from the Riemann sphere to the Riemann sphere."

"These transformations preserve angles, map every straight line to a line or circle, and map every circle to a line or circle."

All information is probided from (Wikipedia, 2020). These exerpts provide the definition of a Mobius transformation as well as an invarient within the transformation.

\subsection{Write down a math term/word/notation that was totally unfamiliar to you.}
Irrational Rotation.

\subsection{Using a resource (wikipedia is fine), provide a formal definition of the word or symbol as it was intended in the presentation.}
In the mathematical theory of dynamicval systems, an irrational rotation is a map where $\theta$ is an irrational number. Under the identification of a circle with $R/z$ or with the interval $[0,1]$ with the boundary points glued together, this map becomes a rotation of a circle by a proportion $\theta$ of a full revolution.

\subsection{Ask a question that the presentation raised for you that you would like to know more about.}
Why isn't any orbit periodic? Within the same slide, Dr. Hamidi makes a statement "No orbit is periodic." What is a non-periodic orbit?

\newpage

\section{Question 4}

\subsection{Provide the full statement of at least one theorem, as it’s stated on the presenter’s slides.}
Every $C^*$-algebra looks like a closed $*$-subalgebra of $B(H)$ for some Hilbert space $H$.

\subsection{In the theorem statement above, what type of a proposition is it? Is it conditional? Does it have suppressed quantifiers? If so, rewrite the theorem statement showing these.}
It is a universal statement because the theorem states that ``Every $C^*$-algebra...'' an example of this is:

\begin{center}
  $c(x) = \{f : x \rightarrow \mathbb{C} : $f is continuous$\}$ where x is compact, Hausdorff...
\end{center}

I suppose that the theorem could be conditional because it has to be a $C^*$-algebra.

We can write this as
\begin{center}
  $\forall C^* algebra~\exists H ~s.t. ~C^* algebra~ = ~closed~*$-subalgebra$~ \subseteq B(H)$.
\end{center}

\subsection{Write down a math term/word/notation that seems familiar to you (look out for words like injective, surjective, or notation like function notation).}
Let $(A(\mathbb{D}), \alpha)$ be a dynamical system. Then there exists a Mobius transformation of the form

\begin{center}
  $\tau(w) = \lambda (\frac{w-\mu}{1-\mu\bar w})$
\end{center}
such that $\alpha_1(f)(w) = f(\tau(w))$.

Within this I understand what it means for a system to be dynamical and for a transformation between spaces to exist.

\subsection{What do you understand this word/notation to mean?}
I understand this exerpt to mean there is some transform called a Mobius transform that exists as a function that accepts an argument $w$ and uses it to convert into something that is possibly called Mobius space.

\subsection{Using a resource (wikipedia is fine), provide a formal definition of the word or symbol as it was intended in the presentation.}
Here I look into the definition of a Mobius transformation.

In geometry and complex analysis, a Möbius transformation of the complex plane is a rational function of the form

\begin{center}
  $f(z) = \frac{az+b}{cz+d}$
\end{center}
of one complex variable $z$; here the coefficients $a,b,c,d$ are complex numbers satisfying $ab-bc \neq 0$.

"A Mobius transformation is always a bijective holomorphic function from the Riemann sphere to the Riemann sphere."

"These transformations preserve angles, map every straight line to a line or circle, and map every circle to a line or circle."

All information is probided from (Wikipedia, 2020). These exerpts provide the definition of a Mobius transformation as well as an invarient within the transformation.

\subsection{Write down a math term/word/notation that was totally unfamiliar to you.}
Irrational Rotation.

\subsection{Using a resource (wikipedia is fine), provide a formal definition of the word or symbol as it was intended in the presentation.}
In the mathematical theory of dynamicval systems, an irrational rotation is a map where $\theta$ is an irrational number. Under the identification of a circle with $R/z$ or with the interval $[0,1]$ with the boundary points glued together, this map becomes a rotation of a circle by a proportion $\theta$ of a full revolution.

\subsection{Ask a question that the presentation raised for you that you would like to know more about.}
Why isn't any orbit periodic? Within the same slide, Dr. Hamidi makes a statement "No orbit is periodic." What is a non-periodic orbit?

\newpage

\section{Question 5}

\subsection{Provide the full statement of at least one theorem, as it’s stated on the presenter’s slides.}
Every $C^*$-algebra looks like a closed $*$-subalgebra of $B(H)$ for some Hilbert space $H$.

\subsection{In the theorem statement above, what type of a proposition is it? Is it conditional? Does it have suppressed quantifiers? If so, rewrite the theorem statement showing these.}
It is a universal statement because the theorem states that ``Every $C^*$-algebra...'' an example of this is:

\begin{center}
  $c(x) = \{f : x \rightarrow \mathbb{C} : $f is continuous$\}$ where x is compact, Hausdorff...
\end{center}

I suppose that the theorem could be conditional because it has to be a $C^*$-algebra.

We can write this as
\begin{center}
  $\forall C^* algebra~\exists H ~s.t. ~C^* algebra~ = ~closed~*$-subalgebra$~ \subseteq B(H)$.
\end{center}

\subsection{Write down a math term/word/notation that seems familiar to you (look out for words like injective, surjective, or notation like function notation).}
Let $(A(\mathbb{D}), \alpha)$ be a dynamical system. Then there exists a Mobius transformation of the form

\begin{center}
  $\tau(w) = \lambda (\frac{w-\mu}{1-\mu\bar w})$
\end{center}
such that $\alpha_1(f)(w) = f(\tau(w))$.

Within this I understand what it means for a system to be dynamical and for a transformation between spaces to exist.

\subsection{What do you understand this word/notation to mean?}
I understand this exerpt to mean there is some transform called a Mobius transform that exists as a function that accepts an argument $w$ and uses it to convert into something that is possibly called Mobius space.

\subsection{Using a resource (wikipedia is fine), provide a formal definition of the word or symbol as it was intended in the presentation.}
Here I look into the definition of a Mobius transformation.

In geometry and complex analysis, a Möbius transformation of the complex plane is a rational function of the form

\begin{center}
  $f(z) = \frac{az+b}{cz+d}$
\end{center}
of one complex variable $z$; here the coefficients $a,b,c,d$ are complex numbers satisfying $ab-bc \neq 0$.

"A Mobius transformation is always a bijective holomorphic function from the Riemann sphere to the Riemann sphere."

"These transformations preserve angles, map every straight line to a line or circle, and map every circle to a line or circle."

All information is probided from (Wikipedia, 2020). These exerpts provide the definition of a Mobius transformation as well as an invarient within the transformation.

\subsection{Write down a math term/word/notation that was totally unfamiliar to you.}
Irrational Rotation.

\subsection{Using a resource (wikipedia is fine), provide a formal definition of the word or symbol as it was intended in the presentation.}
In the mathematical theory of dynamicval systems, an irrational rotation is a map where $\theta$ is an irrational number. Under the identification of a circle with $R/z$ or with the interval $[0,1]$ with the boundary points glued together, this map becomes a rotation of a circle by a proportion $\theta$ of a full revolution.

\subsection{Ask a question that the presentation raised for you that you would like to know more about.}
Why isn't any orbit periodic? Within the same slide, Dr. Hamidi makes a statement "No orbit is periodic." What is a non-periodic orbit?

\newpage

\section{Question 6}
\begin{definition}
The sequence $(a_n)$ diverges as $n \rightarrow \infty$ if and only if for any large number $M > 0$, there exists $N \in \mathbb{N}$ such that $|a(n)| > M$ for all $n \geq N$.
\end{definition}

Define a sequence $a_n := \frac{2n-1}{2}$ for each $n \in \mathbb{Z}^+$.

\subsection{Use induction to prove $\sum_{n=1}^k a_n = \frac{k^2}{2}$}
\subsubsection{Define a predicate $P(k)$.}
\begin{center}
$P(k) := \sum_{n=1}^k a_n = \frac{k^2}{2}$
\end{center}

\subsubsection{State what you are trying to show by quantifying your predicate.}
\begin{center}
$\forall k \in \mathbb{Z}^+ P(k)$
\end{center}

\subsubsection{Show $P(1)$ is true.}
\begin{center}
$P(1) = (\sum_{n=1}^1 a_n = \frac{1^2}{2})$ \\
$P(1) = (\frac{2 - 1}{2} = \frac{1}{2}$) \\
$P(1) = (\frac{1}{2} = \frac{1}{2})$ \\
$P(1)$ is True
\end{center}

\subsubsection{Prove that if $P(m)$ is true for some $k \in \mathbb{Z}^+$, then $P(m+1)$ is true.}
\begin{center}
\begin{tabular}{lll}
$P(m)$ is True. && Given \\

$P(m+1) := (\sum_{n=1}^{m+1} a_{n} = \frac{(m+1)^2}{2})$ & $\iff$ \\
$P(m+1) := (a_{m+1} + \sum_{n=1}^{m} a_{n} = \frac{m^2 + 2m + 1}{2})$ & $\iff$ \\
$P(m+1) := (a_{m+1} + \sum_{n=1}^{m} a_{n} = \frac{m^2}{2} + \frac{2m + 1}{2})$ & $\iff$ \\
$P(m+1) := (a_{m+1} = \frac{2m+1}{2})$ & $\iff$ & [1] \\
$P(m+1) := (\frac{2(m+1) - 1}{2} = \frac{2m + 1}{2})$ & $\iff$ \\
$P(m+1) := (\frac{2m + 1}{2} = \frac{2m + 1}{2})$ \\
$P(m+1)$ is True \\
\end{tabular} \\
$[1]$ Terms cancel because $P(m)$ is True.
\end{center}

This shows directly that we can remove a significant portion of the burden of the claim from $P(m+1)$ by knowing $P(m)$ is true.

\subsection{Prove that the sequence $S_k := \sum_{n=1}^k a_n$ diverges as $k \rightarrow \infty$}
The above proof has demonstrated  that this summation can be represented as the closed form formula $S_k := \frac{k^2}{2}$. This means we can demonstrate that $\frac{k^2}{2}$ diverges to signify that $\sum_{n=1}^k a_n$ diverges.

Consider some arbitrarily large number $M$ s.t. $M > 0$ Does there exist some $N \in \mathbb{N}$ such that $|S_k| > M$ for all $n \geq N$? If this $N$ exists it means that $S_N > M \equiv \frac{N^2}{2} > M$.

Algebraically solving for $N$ we find that $N > \sqrt{2M}$. This shows that there exists some $N$ that causes $S_N$ to be larger than M in all cases. Since $N \in \mathbb{N}$ we can see that $\frac{N^2}{2}$ is a strictly increasing function. This implies that if $S_N > M \rightarrow \forall n \geq N (S_n > M)$. This is the definition of a diverging series. Therefore $S_k$ diverges as $k \rightarrow \infty$.


\end{document}
